\section{Wave-branch structure and characteristic speeds}
\label{sec:waves}

The substrate carries Lorentzian sign structure only at the level of the brane \emph{action}
(\Cref{subsec:ontology-kinematics}); the ambient $\R^4$ is symmetric Euclidean.
This section collects the linearized dispersion results derived in \Cref{sec:continuum-model}
and states them as exported interface equations for companion papers.
The deeper question of whether inside observers built from the same substrate inherit universal
effective Lorentz kinematics is the subject of the Lorentz bridge paper and is not claimed here.

\subsection{Linearized dispersion and branch speeds}

Linearizing \Cref{eq:eom} about a homogeneous background yields a multi-branch wave system.
The characteristic speeds $c_L^2 = k_s a^2/m$ and $c_T^2 = (1-\alpha)\,k_s a^2/m$ are derived
in \Cref{sec:continuum-model} from two independent starting points --- the closed-form dispersion
of the discrete action (\Cref{eq:cone-speeds}) and the acoustic-tensor eigenvalues of the $W_2$
continuum limit (\Cref{eq:wave-speeds-derived}) --- and their agreement there is a self-consistency
check across levels of description.
For each polarization branch $b$ one has the long-wavelength dispersion
\begin{equation}
  \omega_b^2(\veck)\;\approx\; c_b^2\,|\veck|^2 \quad \text{for } |\veck|\to 0,
  \label{eq:linear-dispersion}
\end{equation}
which defines an effective wave cone for that branch.

Two structural points follow immediately.
First, the prestress parameter $\alpha$ controls the \emph{splitting} between transverse and
longitudinal speeds: $\alpha=0$ collapses them to a single speed $c_L$, while $\alpha\to 1$
drives $c_T\to 0$, approaching a degenerate shear-free medium.
Second, the transverse branch carries the geometric coupling
$\partial_i u\,\partial_j u$ as its leading nonlinearity (\Cref{sec:continuum-model}), which
is absent at linear order but feeds back at finite amplitude --- the mechanism underlying
the gravity-like contraction channel (\Cref{subsec:gravity-channel}).

\subsection{Regime of validity}
\label{subsec:wave-regime}

The single-branch description \eqref{eq:linear-dispersion} is valid when:
(i) a single branch dominates transport and mode mixing is negligible,
(ii) dispersion at shorter wavelengths is small, and
(iii) geometric nonlinearity at larger amplitudes is not yet important.
In the present theory, departures from these conditions are not treated as flaws;
they are structural consequences of a substrate whose timelike direction is selected by
the brane action rather than imposed as a metric primitive, and they delimit the domain
where a single-branch or relativistic effective description is expected to apply.

\paragraph{Interface to the Lorentz bridge paper.}
The Lorentz bridge paper imports \Cref{eq:linear-dispersion,eq:wave-speeds-derived}
and the dual-observer argument (\Cref{subsec:gravity-channel}) to ask whether inside
observers built from the same substrate inherit universal effective Lorentz kinematics
at the wave cone defined by $c_T$.
Residual birefringence ($c_L \neq c_T$) is then the primary falsifiable signature.